% ================================================================
%  SESIÓN 08 — Área de zonas urbanas
%  Almería en Cifras y Formas · Matemáticas 1.º ESO
%  IES Al-Ándalus · 2025-2026
% ================================================================
\documentclass[aspectratio=169, 10pt, spanish]{beamer}
\usepackage{almeria-beamer}

\renewcommand{\SessionNumber}{08}
\renewcommand{\SessionTitle}{Área de zonas urbanas}
\renewcommand{\SessionName}{Sesión 08}
\renewcommand{\SessionObjective}{%
  Calcular el área de figuras planas básicas y aplicarla a zonas verdes
  y plazas de Almería, interpretando correctamente las unidades
  de superficie.%
}
\renewcommand{\BlockName}{Bloque 2: Geometría urbana}

\begin{document}

% ---------------------------------------------------------------
%  PORTADA
% ---------------------------------------------------------------
\begin{frame}[plain]
  \titlepage
\end{frame}

% ---------------------------------------------------------------
%  ÍNDICE
% ---------------------------------------------------------------
\begin{frame}{Contenidos de hoy}
  \begin{columns}[t]
    \begin{column}{0.48\textwidth}
      \begin{itemize}
        \item Definición de área
        \item Analogía: el césped del parque
        \item Unidades de superficie
        \item Fórmulas de las 6 figuras básicas
        \item Método de cuadrícula
      \end{itemize}
    \end{column}
    \begin{column}{0.48\textwidth}
      \begin{itemize}
        \item Conversiones de unidades
        \item Ejemplos en Almería
        \item Problema paso a paso
        \item Ejercicios Bloom
        \item Evidencia del día
      \end{itemize}
    \end{column}
  \end{columns}
  \vspace{0.4cm}
  \begin{notabox}[Niveles de Bloom]
    \bloomtag{Recordar}\quad
    \bloomtag{Comprender}\quad
    \bloomtag{Aplicar}
  \end{notabox}
\end{frame}

% ---------------------------------------------------------------
\seccionframe{¿Qué es el área?}
% ---------------------------------------------------------------

% ---------------------------------------------------------------
%  FRAME 3 — Definición
% ---------------------------------------------------------------
\begin{frame}{Área: la superficie interior}
  \begin{columns}[c]
    \begin{column}{0.54\textwidth}
      \begin{definicionbox}{Área}
        El \textbf{área} de una figura es la medida de la
        \textbf{superficie encerrada} en su interior.
        Indica cuánto espacio ocupa la figura en el plano.
      \end{definicionbox}

      \vspace{0.35cm}
      \begin{notabox}[Analogía]
        \AlmFontSmall El área es como el \textbf{césped} a segar dentro
        del parque: mides cuánto verde hay,
        no la longitud de la valla que lo rodea.
      \end{notabox}

      \vspace{0.3cm}
      \textbf{Unidades:} m$^2$, km$^2$, cm$^2$, mm$^2$, ha, a \\
      \AlmFontSmall (dos dimensiones $\Rightarrow$ el exponente es 2)
    \end{column}
    \begin{column}{0.42\textwidth}
      \begin{center}
      \begin{tikzpicture}[scale=0.80]
        % Parque con cuadrícula de área
        \draw[AlmAzulC!30, very thin, step=0.5]
          (0,0) grid (3.5,2.5);
        \draw[AlmTerra, line width=2.5pt] (0,0) rectangle (3.5,2.5);
        \fill[AlmVerde!22] (0,0) rectangle (3.5,2.5);
        \draw[AlmAzulC!30, very thin, step=0.5]
          (0,0) grid (3.5,2.5);
        \node[font=\AlmFontSmall\bfseries, color=AlmVerde!70!black]
          at (1.75,1.25) {ÁREA};
        \node[font=\AlmFontTiny\bfseries, color=AlmTerra]
          at (1.75,-0.3) {Perímetro (borde)};
        % Cota
        \draw[acota] (0,-0.65) -- (3.5,-0.65)
          node[midway,below,font=\AlmFontTiny]{base $b$};
        \draw[acota] (3.9,0) -- (3.9,2.5)
          node[midway,right,font=\AlmFontTiny]{altura $h$};
      \end{tikzpicture}
      \end{center}
    \end{column}
  \end{columns}
\end{frame}

% ---------------------------------------------------------------
%  FRAME 4 — Unidades de superficie
% ---------------------------------------------------------------
\begin{frame}{Unidades de superficie y conversiones}
  \begin{columns}[c]
    \begin{column}{0.54\textwidth}
      \begin{center}
      \AlmFontSmall
      \begin{tabular}{>{\bfseries\color{AlmAzulM}}l l l}
        \toprule
        Unidad & Símbolo & Equivalencia \\
        \midrule
        milímetro cuadrado & mm$^2$ & \\
        centímetro cuadrado & cm$^2$ & 1 cm$^2$ = 100 mm$^2$ \\
        metro cuadrado & m$^2$ & 1 m$^2$ = 10{,}000 cm$^2$ \\
        área & a & 1 a = 100 m$^2$ \\
        hectárea & ha & 1 ha = 10{,}000 m$^2$ \\
        kilómetro cuadrado & km$^2$ & 1 km$^2$ = 1{,}000{,}000 m$^2$ \\
        \bottomrule
      \end{tabular}
      \end{center}
    \end{column}
    \begin{column}{0.42\textwidth}
      \begin{formulabox}[AlmAmbar!60!black]
        \AlmFontSmall
        1 ha = 10{,}000 m$^2$\\[4pt]
        1 km$^2$ = 100 ha
      \end{formulabox}
      \vspace{0.3cm}
      \begin{notabox}[Uso práctico]
        \AlmFontSmall
        \begin{itemize}
          \item Parcelas/campos: \textbf{ha}
          \item Pisos/plazas: \textbf{m$^2$}
          \item Comunidades/países: \textbf{km$^2$}
        \end{itemize}
      \end{notabox}
    \end{column}
  \end{columns}
\end{frame}

% ---------------------------------------------------------------
%  FRAME 5 — Fórmulas: cuadrado y rectángulo
% ---------------------------------------------------------------
\begin{frame}{Fórmulas de área (I): cuadrado y rectángulo}
  \begin{columns}[c]
    \begin{column}{0.50\textwidth}
      \begin{center}
      \begin{tikzpicture}[scale=0.80]
        \fill[AlmVerde!15] (0,0) rectangle (2.0,2.0);
        \draw[rectov, line width=1.5pt] (0,0) rectangle (2.0,2.0);
        \draw[acota] (0,-0.5) -- (2.0,-0.5)
          node[midway,below,font=\AlmFontScript]{$l$};
        \draw[acota] (2.4,0) -- (2.4,2.0)
          node[midway,right,font=\AlmFontScript]{$l$};
        \node[font=\AlmFontSmall\bfseries, color=AlmVerde!70!black]
          at (1.0,1.0) {$A = l^2$};
        \node[below, font=\AlmFontScript\bfseries, color=AlmVerde]
          at (1.0,-0.9) {Cuadrado};
      \end{tikzpicture}
      \end{center}
      \formuladest[AlmVerde]{A_{\text{cuad}} = l^2}
    \end{column}
    \begin{column}{0.46\textwidth}
      \begin{center}
      \begin{tikzpicture}[scale=0.80]
        \fill[AlmAzulC!18] (0,0) rectangle (2.8,1.6);
        \draw[zona, line width=1.5pt] (0,0) rectangle (2.8,1.6);
        \draw[acota] (0,-0.5) -- (2.8,-0.5)
          node[midway,below,font=\AlmFontScript]{$b$};
        \draw[acota] (3.2,0) -- (3.2,1.6)
          node[midway,right,font=\AlmFontScript]{$h$};
        \node[font=\AlmFontSmall\bfseries, color=AlmAzulM]
          at (1.4,0.8) {$A = b \times h$};
        \node[below, font=\AlmFontScript\bfseries, color=AlmAzulM]
          at (1.4,-0.9) {Rectángulo};
      \end{tikzpicture}
      \end{center}
      \formuladest[AlmAzulM]{A_{\text{rect}} = b \times h}
    \end{column}
  \end{columns}
\end{frame}

% ---------------------------------------------------------------
%  FRAME 6 — Fórmulas: triángulo y paralelogramo
% ---------------------------------------------------------------
\begin{frame}{Fórmulas de área (II): triángulo y paralelogramo}
  \begin{columns}[c]
    \begin{column}{0.50\textwidth}
      \begin{center}
      \begin{tikzpicture}[scale=0.80]
        \fill[AlmTerra!12] (0,0) -- (3.0,0) -- (2.0,1.8) -- cycle;
        \draw[rectot, line width=1.5pt] (0,0) -- (3.0,0) -- (2.0,1.8) -- cycle;
        % Altura
        \draw[AlmGris, dashed, thin] (2.0,0) -- (2.0,1.8);
        \draw[acota] (-0.4,0) -- (-0.4,1.8)
          node[midway,left,font=\AlmFontScript]{$h$};
        \draw[acota] (0,-0.5) -- (3.0,-0.5)
          node[midway,below,font=\AlmFontScript]{$b$};
        \node[font=\AlmFontSmall\bfseries, color=AlmTerra]
          at (1.5,0.6) {$\dfrac{b \times h}{2}$};
        \node[below, font=\AlmFontScript\bfseries, color=AlmTerra]
          at (1.5,-0.9) {Triángulo};
      \end{tikzpicture}
      \end{center}
      \formuladest[AlmTerra]{A_{\text{trián}} = \dfrac{b \times h}{2}}
    \end{column}
    \begin{column}{0.46\textwidth}
      \begin{center}
      \begin{tikzpicture}[scale=0.80]
        \fill[BAnalizar!10]
          (0.5,0) -- (3.1,0) -- (2.6,1.6) -- (0.0,1.6) -- cycle;
        \draw[rectov, line width=1.5pt]
          (0.5,0) -- (3.1,0) -- (2.6,1.6) -- (0.0,1.6) -- cycle;
        \draw[AlmGris, dashed, thin] (0.5,0) -- (0.5,1.6);
        \draw[acota] (-0.4,0) -- (-0.4,1.6)
          node[midway,left,font=\AlmFontScript]{$h$};
        \draw[acota] (0.5,-0.5) -- (3.1,-0.5)
          node[midway,below,font=\AlmFontScript]{$b$};
        \node[font=\AlmFontSmall\bfseries, color=AlmVerde]
          at (1.6,0.8) {$A = b \times h$};
        \node[below, font=\AlmFontScript\bfseries, color=AlmVerde]
          at (1.6,-0.9) {Paralelogramo};
      \end{tikzpicture}
      \end{center}
      \formuladest[BAnalizar]{A_{\text{paral}} = b \times h}
    \end{column}
  \end{columns}
\end{frame}

% ---------------------------------------------------------------
%  FRAME 7 — Fórmulas: trapecio y círculo
% ---------------------------------------------------------------
\begin{frame}{Fórmulas de área (III): trapecio y círculo}
  \begin{columns}[c]
    \begin{column}{0.50\textwidth}
      \begin{center}
      \begin{tikzpicture}[scale=0.80]
        \fill[AlmAmbar!15]
          (0.4,0) -- (3.2,0) -- (2.8,1.6) -- (0.8,1.6) -- cycle;
        \draw[AlmAmbar!70!black, line width=1.5pt]
          (0.4,0) -- (3.2,0) -- (2.8,1.6) -- (0.8,1.6) -- cycle;
        \draw[acota] (0.4,-0.5) -- (3.2,-0.5)
          node[midway,below,font=\AlmFontTiny]{$B$ (base mayor)};
        \draw[acota] (0.8,1.9) -- (2.8,1.9)
          node[midway,above,font=\AlmFontTiny]{$b$ (base menor)};
        \draw[AlmGris, dashed, thin] (0.8,0) -- (0.8,1.6);
        \draw[acota] (-0.35,0) -- (-0.35,1.6)
          node[midway,left,font=\AlmFontTiny]{$h$};
        \node[below, font=\AlmFontScript\bfseries, color=AlmAmbar!80!black]
          at (1.8,-0.9) {Trapecio};
      \end{tikzpicture}
      \end{center}
      \formuladest[AlmAmbar!70!black]{A_{\text{trap}} =
        \dfrac{(B+b)}{2} \times h}
    \end{column}
    \begin{column}{0.46\textwidth}
      \begin{center}
      \begin{tikzpicture}[scale=0.80]
        \fill[AlmAzulC!20] (1.5,1.5) circle (1.4);
        \draw[AlmAzulM, line width=1.5pt] (1.5,1.5) circle (1.4);
        \draw[AlmTerra, line width=1pt] (1.5,1.5) -- (2.9,1.5)
          node[midway,above,font=\AlmFontScript]{$r$};
        \node[font=\AlmFontSmall\bfseries, color=AlmAzulM]
          at (1.5,1.5) {$\pi r^2$};
        \node[below, font=\AlmFontScript\bfseries, color=AlmAzulM]
          at (1.5,0.0) {Círculo};
      \end{tikzpicture}
      \end{center}
      \formuladest[AlmAzulM]{A_{\text{círc}} = \pi \times r^2}
    \end{column}
  \end{columns}
\end{frame}

% ---------------------------------------------------------------
\seccionframe[BAplicar]{Áreas en Almería}
% ---------------------------------------------------------------

% ---------------------------------------------------------------
%  FRAME 8 — Ejemplos en Almería
% ---------------------------------------------------------------
\begin{frame}{Áreas de zonas verdes de Almería}
  \begin{columns}[t]
    \begin{column}{0.54\textwidth}
      \begin{ejemplobox}
        \AlmFontSmall
        \textbf{Parque Nicolás Salmerón}
        (rectángulo 400 m $\times$ 80 m):
        \[A = 400 \times 80 = \mathbf{32{,}000\ m^2} = \mathbf{3{,}2\ ha}\]

        \vspace{0.1cm}
        \textbf{Plaza triangular} (base 120 m, altura 95 m):
        \[A = \frac{120 \times 95}{2} = \frac{11{,}400}{2}
          = \mathbf{5{,}700\ m^2}\]

        \vspace{0.1cm}
        \textbf{Alcazaba de Almería}
        (área aproximada):
        \[A \approx \mathbf{43{,}000\ m^2} = \mathbf{4{,}3\ ha}\]
      \end{ejemplobox}
    \end{column}
    \begin{column}{0.42\textwidth}
      \begin{notabox}[¿Qué es una hectárea?]
        \AlmFontSmall
        1 ha = 10{,}000 m$^2$.
        \\[4pt]
        Equivale a un cuadrado de
        \textbf{100 m} $\times$ \textbf{100 m}.
        \\[4pt]
        Se usa en agricultura, urbanismo
        y cartografía.
        \\[4pt]
        La Alcazaba ocupa más de 4 campos
        de fútbol (un campo $\approx$ 0{,}7 ha).
      \end{notabox}
    \end{column}
  \end{columns}
\end{frame}

% ---------------------------------------------------------------
%  FRAME 9 — Problema paso a paso
% ---------------------------------------------------------------
\begin{frame}{Problema resuelto: manzana trapezoidal}
  \begin{definicionbox}{Enunciado}
    \AlmFontSmall Una manzana de Almería tiene forma de trapecio con
    bases paralelas de \textbf{200 m} y \textbf{140 m},
    y altura \textbf{90 m}.
    Calcula su área en m$^2$ y en hectáreas.
  \end{definicionbox}

  \vspace{0.15cm}
  \begin{columns}[t]
    \begin{column}{0.50\textwidth}
      \begin{pasobox}[1]
        \AlmFontSmall Identificar la fórmula del trapecio:
        \[A = \frac{(B + b)}{2} \times h\]
      \end{pasobox}
      \vspace{0.2cm}
      \begin{pasobox}[2]
        \AlmFontSmall Sustituir valores:
        \[A = \frac{(200 + 140)}{2} \times 90
          = \frac{340}{2} \times 90\]
      \end{pasobox}
    \end{column}
    \begin{column}{0.46\textwidth}
      \begin{pasobox}[3]
        \AlmFontSmall Calcular:
        \[A = 170 \times 90 = \mathbf{15{,}300\ m^2}\]
      \end{pasobox}
      \vspace{0.2cm}
      \begin{pasobox}[4]
        \AlmFontSmall Convertir a hectáreas:
        \[A = \frac{15{,}300}{10{,}000} = \mathbf{1{,}53\ ha}\]
      \end{pasobox}
    \end{column}
  \end{columns}
\end{frame}

% ---------------------------------------------------------------
\seccionframe[BRecordar]{Ejercicios — RECORDAR}
% ---------------------------------------------------------------

% ---------------------------------------------------------------
%  FRAME 10 — Ejercicio RECORDAR
% ---------------------------------------------------------------
\begin{frame}{Ejercicio 1 — RECORDAR \dificultad{1}}
  \begin{recordarbox}
    \begin{enumerate}
      \item Escribe la fórmula del área de:
            rectángulo, triángulo y círculo.
      \item Relaciona cada figura con su fórmula:

      \vspace{0.2cm}
      \begin{center}
      \AlmFontSmall
      \begin{tabular}{|l|l|}
        \hline
        \rowcolor{AlmAzul!15}
        \textbf{Figura} & \textbf{Fórmula} \\
        \hline
        Cuadrado & $A = \pi r^2$ \\
        \hline
        Triángulo & $A = l^2$ \\
        \hline
        Trapecio & $A = b \times h / 2$ \\
        \hline
        Círculo & $A = \tfrac{(B+b)}{2} \times h$ \\
        \hline
      \end{tabular}
      \end{center}
    \end{enumerate}
  \end{recordarbox}
  \vspace{0.15cm}
  \begin{notabox}[Pista]
    \AlmFontSmall Para el triángulo, la altura debe ser perpendicular a la base.
  \end{notabox}
\end{frame}

% ---------------------------------------------------------------
\seccionframe[BComprender]{Ejercicios — COMPRENDER}
% ---------------------------------------------------------------

% ---------------------------------------------------------------
%  FRAME 11 — Ejercicio COMPRENDER
% ---------------------------------------------------------------
\begin{frame}{Ejercicio 2 — COMPRENDER \dificultad{2}}
  \begin{comprenderbox}
    Un \textbf{cuadrado} de lado 10 m y un \textbf{rectángulo}
    de 15 m $\times$ 5 m tienen el mismo perímetro (40 m).
    \begin{enumerate}
      \item Calcula el área de cada uno.
      \item ¿Cuál tiene mayor área?
      \item ¿Qué conclusión extraes sobre la relación
            entre perímetro y área?
    \end{enumerate}
  \end{comprenderbox}

  \vspace{0.2cm}
  \begin{columns}[c]
    \begin{column}{0.40\textwidth}
      \begin{center}
      \begin{tikzpicture}[scale=0.60]
        \fill[AlmVerde!15] (0,0) rectangle (2.0,2.0);
        \draw[rectov] (0,0) rectangle (2.0,2.0);
        \draw[acota] (0,-0.5) -- (2.0,-0.5)
          node[midway,below,font=\AlmFontTiny]{10 m};
        \node[font=\AlmFontScript\bfseries, color=AlmVerde] at (1.0,1.0) {?};
      \end{tikzpicture}
      \end{center}
    \end{column}
    \begin{column}{0.55\textwidth}
      \begin{center}
      \begin{tikzpicture}[scale=0.60]
        \fill[AlmAzulC!18] (0,0) rectangle (3.0,1.0);
        \draw[zona] (0,0) rectangle (3.0,1.0);
        \draw[acota] (0,-0.5) -- (3.0,-0.5)
          node[midway,below,font=\AlmFontTiny]{15 m};
        \draw[acota] (3.4,0) -- (3.4,1.0)
          node[midway,right,font=\AlmFontTiny]{5 m};
        \node[font=\AlmFontScript\bfseries, color=AlmAzulM] at (1.5,0.5) {?};
      \end{tikzpicture}
      \end{center}
    \end{column}
  \end{columns}
\end{frame}

% ---------------------------------------------------------------
\seccionframe[BAplicar]{Ejercicios — APLICAR}
% ---------------------------------------------------------------

% ---------------------------------------------------------------
%  FRAME 12 — Ejercicio APLICAR
% ---------------------------------------------------------------
\begin{frame}{Ejercicio 3 — APLICAR \dificultad{3}}
  \begin{aplicarbox}
    \AlmFontSmall Se quiere crear una \textbf{zona verde triangular}
    en la playa de Almería con base \textbf{300 m}
    y altura \textbf{150 m}.
    \begin{enumerate}
      \item Calcula el área en m$^2$.
      \item Exprésala en hectáreas (ha).
      \item Calcula el coste de sembrar césped
            si el precio es \textbf{8 €/m$^2$}.
    \end{enumerate}
  \end{aplicarbox} 
  \begin{columns}[c]
    \begin{column}{0.42\textwidth}
      \begin{center}
      \begin{tikzpicture}[scale=0.55]
        \fill[AlmVerde!18] (0,0) -- (4.5,0) -- (2.25,2.25) -- cycle;
        \draw[rectov, line width=1.5pt]
          (0,0) -- (4.5,0) -- (2.25,2.25) -- cycle;
        \draw[AlmGris, dashed, thin] (2.25,0) -- (2.25,2.25);
        \draw[acota] (0,-0.6) -- (4.5,-0.6)
          node[midway,below,font=\AlmFontTiny]{300 m};
        \draw[acota] (2.6,0) -- (2.6,2.25)
          node[midway,right,font=\AlmFontTiny]{150 m};
        \node[font=\AlmFontScript\bfseries, color=AlmVerde!70!black]
          at (2.25,0.8) {Zona verde};
      \end{tikzpicture}
      \end{center}
      
            \begin{pasobox}[1]
              \AlmFontScript $A = \dfrac{b \times h}{2}
                = \dfrac{300 \times 150}{2} = ?\ \text{m}^2$
            \end{pasobox}
    \end{column}
    \begin{column}{0.54\textwidth}
      \begin{pasobox}[2]
        \AlmFontScript $A_{\text{ha}} = \dfrac{A_{\text{m}^2}}{10{,}000}$
      \end{pasobox}
      \vspace{0.15cm}
      \begin{pasobox}[3]
        \AlmFontScript $\text{Coste} = A_{\text{m}^2} \times 8$~€/m$^2$
      \end{pasobox}
    \end{column}
  \end{columns}
\end{frame}

% ---------------------------------------------------------------
%  FRAME 13 — Error frecuente
% ---------------------------------------------------------------
\begin{frame}{¡Cuidado con este error!}
  \begin{errorbox}
    \begin{itemize}
      \item \incorrecto\
            \textbf{«El parque tiene 32{,}000 metros»}
            \\[6pt]
            \correcto\
            El área se mide en metros \textbf{cuadrados}:
            «El parque tiene 32{,}000 \textbf{m$^2$}» o «3{,}2 ha».
            \\
            La palabra «metros» sola indica longitud (perímetro, distancia),
            no superficie.

      \vspace{0.4cm}
      \item \incorrecto\
            Usar la fórmula del rectángulo en un triángulo:
            \\
            $A = 120 \times 95 = 11{,}400$~m$^2$ (sin dividir entre 2).
            \\[6pt]
            \correcto\
            $A = \dfrac{120 \times 95}{2} = 5{,}700$~m$^2$
            \\
            El triángulo es \textbf{la mitad} del rectángulo que lo contiene.
    \end{itemize}
  \end{errorbox}
\end{frame}

% ---------------------------------------------------------------
%  FRAME 14 — Resumen
% ---------------------------------------------------------------
\begin{frame}{Resumen de la sesión}
  \begin{resumenbox}
    \begin{itemize}
      \item El \textbf{área} mide la superficie interior:
            como el césped a segar.
      \item Unidades: m$^2$, cm$^2$, km$^2$, \textbf{ha}
            (1 ha = 10{,}000 m$^2$).
      \item Fórmulas clave:
        $A_{\text{cuad}} = l^2$, \quad
        $A_{\text{rect}} = b \times h$, \quad
        $A_{\text{trián}} = \tfrac{bh}{2}$, \\
        $A_{\text{paral}} = bh$, \quad
        $A_{\text{trap}} = \tfrac{B+b}{2} \times h$, \quad
        $A_{\text{círc}} = \pi r^2$.
      \item La altura debe ser \textbf{perpendicular} a la base.
      \item Parque Salmerón: $A = 32{,}000$~m$^2 = 3{,}2$~ha.
    \end{itemize}
  \end{resumenbox}
\end{frame}

% ---------------------------------------------------------------
%  FRAME 15 — Evidencia del día
% ---------------------------------------------------------------
\begin{frame}{Evidencia del día}
  \begin{evidenciabox}
    \textbf{Mapa de zonas verdes con áreas calculadas}

    \vspace{0.2cm}
    Investiga \textbf{3 parques o plazas} de Almería.
    Para cada uno:
    \begin{enumerate}
      \item Dibuja su forma aproximada con medidas.
      \item Identifica qué figura geométrica modela mejor su planta.
      \item Aplica la fórmula correcta y calcula el área en m$^2$.
      \item Convierte a hectáreas.
      \item Calcula el coste hipotético de regar el parque
            a razón de 2~€/m$^2$.
    \end{enumerate}
    \vspace{0.2cm}
    \textbf{Criterio de éxito:} los 3 cálculos son correctos,
    las fórmulas están indicadas y las unidades son adecuadas.
  \end{evidenciabox}
\end{frame}

\end{document}



